% Compilen desde esta carpeta o desde la raíz del repositorio.
\makeatletter
\def\input@path{{../../preamble/}{book2026-2027/preamble/}}
\makeatother
\input{preamble.tex}

\usepackage[spanish]{babel}
\usepackage{lmodern}
\usepackage{titlesec}
\geometry{a4paper,margin=1.6cm}
\linespread{1.10}

\definecolor{HeadingColor}{HTML}{243B53}
\definecolor{InstructionColor}{HTML}{704214}
\definecolor{ConceptColor}{HTML}{205C50}
\definecolor{LinkColor}{HTML}{0047AB}
\hypersetup{allcolors=LinkColor}
\newcommand{\guidelink}[2]{\href{#1}{\textbf{\underline{#2}}}}
\newcommand{\instruction}[1]{\textcolor{InstructionColor}{\textbf{#1}}}
\newcommand{\concept}[1]{\textcolor{ConceptColor}{\textbf{#1}}}
\newcommand{\checkbox}{%
    \begingroup
    \setlength{\fboxrule}{0.4pt}%
    \setlength{\fboxsep}{2pt}%
    \fbox{\rule{0pt}{0.8ex}\rule{0.8ex}{0pt}}%
    \endgroup}

\setlength{\parskip}{4pt}
\setlength{\tabcolsep}{5pt}
\renewcommand{\arraystretch}{1.15}
\setlist{leftmargin=0.65cm,topsep=4pt,partopsep=0pt,parsep=0pt,itemsep=3pt}
\titleformat{\section}{\large\bfseries\color{HeadingColor}}{}{0pt}{}
\titlespacing*{\section}{0pt}{10pt}{3pt}
\newcolumntype{Y}{>{\raggedright\arraybackslash}X}

\hypersetup{pdftitle={C5--C6. Apliquen principios de diseño y actualicen el plan},pdfsubject={Emprendimiento e Innovación, Grado 10, Periodo 2}}

\begin{document}

\begin{center}
{\Large\bfseries\color{HeadingColor} C5--C6. Apliquen principios de diseño y actualicen el plan\par}
\vspace{3pt}
Emprendimiento e Innovación \enspace|\enspace Grado 10 \enspace|\enspace Periodo 2\\
Guía 3 de 5
\end{center}

\section*{Objetivo y punto de partida}
\concept{C5.} Aplica principios básicos de diseño, funcionalidad, accesibilidad o experiencia de usuario según el proyecto.\\
\concept{C6.} Organiza recursos, roles y tiempos para avanzar en la construcción de la propuesta.

Recuperen R1, V1 y PL1 (P2, C3--C4), además de D1 y los H pertinentes (P1). Recorran la interacción y localicen qué podría dificultar que el usuario logre su objetivo.

\section*{Documento acumulativo e integridad de la evidencia}
Esta guía es de consulta. Continúen el \textbf{mismo documento del proyecto en LaTeX} en \guidelink{https://www.overleaf.com/}{Overleaf} o en su editor habitual. Actualicen la fecha y entreguen el PDF completo. Usen \textbf{P1} para el Periodo 1 y \textbf{P2} para el Periodo 2 al citar criterios. Enlacen los registros por código o referencia interna y conserven sus fuentes en la sección de origen.

Distingan \concept{evidencia} (lo observado o expresado), \concept{conclusión} (lo que infieren), \concept{supuesto} (lo pendiente de verificar) e \concept{idea} (una propuesta). Registren fuentes y fechas reales. Identifiquen toda simulación autorizada por el docente. \instruction{Si usan IA, limítenla al código LaTeX, al formato y a organizar contenido propio.} Evidencias, conclusiones, ideas, decisiones y aportes deben provenir de su trabajo.

\section*{C5 de P2. Conviertan los principios en decisiones visibles}
\begin{enumerate}[series=pasos]
\item \textbf{Elijan principios pertinentes para su proyecto.} Consideren claridad de instrucciones, facilidad de uso, accesibilidad o respuesta ante una acción. Identifiquen qué dificultad aborda cada principio y relaciónenla con un H o con el recorrido interno de V1. Señalen qué sigue siendo un supuesto.
\item \textbf{Apliquen los principios en el prototipo.} Preparen V2 a partir de V1 y anoten en una foto, captura o secuencia dónde se ve cada decisión. Relacionen el cambio con el paso de R1 correspondiente. Indiquen qué señal observable buscarán después para saber si ayudó.
\end{enumerate}

\textbf{Ejemplo ficticio.} La fila ilustra una decisión y una señal que todavía está por observar.

{\small
\begin{tabularx}{\linewidth}{|Y|Y|Y|}
\hline
\textbf{Principio y dificultad} & \textbf{Decisión y ubicación en V2} & \textbf{Señal para la prueba} \\
\hline
Claridad de la señalización. & Junto al mostrador, la tarjeta indica con texto y una flecha dónde recoger el pedido. & La persona se dirige al puesto correspondiente sin pedir indicaciones. \\
\hline
\end{tabularx}\par
}

\concept{Evidencia de C5.} V2 con cambios localizables y anotaciones que vinculan cada decisión con un principio, una dificultad y una señal para la prueba.

\newpage
\section*{C6 de P2. Actualicen PL1 para la prueba y la siguiente versión}
\begin{enumerate}[resume=pasos]
\item \textbf{Revisen el plan con el equipo.} Vuelvan a PL1 y marquen qué tareas se completaron, cuáles cambiaron y qué quedó pendiente. Registren la fecha de actualización y el motivo de los cambios relevantes. Conserven el plan inicial como referencia dentro del mismo registro.
\item \textbf{Asignen responsables, recursos y plazos.} Añadan las tareas necesarias para dejar V2 lista, organizar la prueba y preparar la siguiente versión. Confirmen quién contactará a los participantes, quién facilitará y quién registrará. Revisen materiales, espacio o acceso digital, permisos y tiempos. Si una tarea depende de otra, indiquen esa dependencia y cómo actuarán si se retrasa.
\end{enumerate}

{\small
\begin{tabularx}{\linewidth}{|Y|Y|Y|Y|Y|}
\hline
\textbf{Tarea o entrega} & \textbf{Responsable y rol} & \textbf{Recursos y acceso} & \textbf{Plazo y estado} & \textbf{Pendiente o dependencia} \\
\hline
Parte de V2, prueba o revisión siguiente. & Quién responde y qué coordina. & Qué hace falta y qué está disponible. & Fecha y avance a la fecha de revisión. & Qué debe confirmarse y quién hará el seguimiento. \\
\hline
\end{tabularx}\par
}

\begin{enumerate}[resume=pasos]
\item \textbf{Comprueben que el plan se puede ejecutar.} Revisen que los tiempos coincidan con la disponibilidad de personas y recursos. Vinculen las tareas al registro de V2 o a la prueba prevista. Actualicen la bitácora con los aportes, los compromisos vigentes y los acuerdos de la sesión.
\end{enumerate}

\concept{Evidencia de C6.} PL1 actualizado y fechado, con responsables, recursos, plazos, estados y pendientes para la prueba y la revisión posterior.

\section*{Resultado y continuidad}
Agreguen V2 anotada y actualicen PL1 en el documento acumulativo. En P2, C7--C8 usarán la versión y las señales de C5 para definir una tarea observable. Consultarán PL1 al realizar la prueba y organizar los ajustes.

\section*{Verificación final}
\begin{itemize}[label=\checkbox]
\item Podemos señalar dónde se refleja cada principio de diseño en V2.
\item Cada cambio responde a una dificultad y tiene una señal observable para la prueba.
\item PL1 conserva la planificación inicial y aclara quién hará qué, con qué y cuándo.
\item Actualizamos la bitácora y compilamos el PDF acumulativo con sus referencias.
\end{itemize}

\end{document}
